% !TeX root = ../教材题库_全部条目_转写工作稿.tex
% Source: XJ-B1-C06-RV-P269-Y0453
% 第6章 / 章末复习 复习题六 / 教材 p.269 / PDF 第276页
% OCR未校对：公式与文字均可能有误；未完成数学TeX转写；禁止直接发给学生
% 原题及相邻原解截图（核对用）：../题目原图/XJ-B1-C06-RV-P269-Y0453-01.webp
% 本文件只在转写完成后作为题干使用；OCR 条目尚不是可靠题干。
\begin{bankitem}{复习题六 第17题}{第6章 / 章末复习 复习题六 / 教材 p.269 / PDF 第276页}{XJ-B1-C06-RV-P269-Y0453}{OCR未校对：公式与文字均可能有误；未完成数学TeX转写；禁止直接发给学生}
\begin{ocrdraft}
[教材p.269 / PDF第276页]
\par
（数学探究活动）在工业产品设计、建筑设计等领域，通常需要参考人体
\par
尺寸数据，以开发符合人体工程学的舒适产品．百分位数在其中发挥了重要的作
\par
用．工程设计师根据国家颁布的成年人人体尺寸设计产品，这将使得更多人适用
\par
该产品．
\par
下表为中国男性人体主要尺寸，节选自《中国成年人人体尺寸》
\par
年龄分组
\par
18、60岁
\par
百
\par
分
\par
位
\par
5
\par
50
\par
99
\par
测量项目
\par
4．1.
\par
1
\par
身高/mm
\par
1543
\par
1583
\par
160
\par
1678
\par
1754
\par
1775
\par
1814
\par
体重/kg
\par
4．1.
\par
2
\par
44
\par
48
\par
50
\par
59
\par
71
\par
4．1.3
\par
上臂长/mm
\par
279
\par
289
\par
294
\par
313
\par
333
\par
338
\par
39
\par
4．1.
\par
4
\par
前臂长/mm
\par
206
\par
216
\par
220
\par
237
\par
253
\par
258
\par
268
\par
4．1.5
\par
大腿长/mm
\par
13
\par
428
\par
436
\par
465
\par
496
\par
505
\par
523
\par
4．1.6
\par
小腿长／mm
\par
324
\par
338
\par
344
\par
369
\par
396
\par
403
\par
419
\par
收集涉及中国成年人人体尺寸的国家标准，了解人体尺寸百分位数的设
\par
定依据．
\par
（2）调查本地区常使用的工业产品数据，如大巴车轿厢高度、公共场所座椅
\par
高度、衣服的尺寸等数据，针对国家标准进行统计探究活动，这些产品符合国家
\par
标准吗？符合个人的舒适感受吗？
\end{ocrdraft}
\end{bankitem}
